%%%%%%%%%%%%%%%%%
% This is an sample CV template created using altacv.cls
% (v1.1.4, 27 July 2018) written by LianTze Lim (liantze@gmail.com). Now compiles with pdfLaTeX, XeLaTeX and LuaLaTeX.
% 
%% It may be distributed and/or modified under the
%% conditions of the LaTeX Project Public License, either version 1.3
%% of this license or (at your option) any later version.
%% The latest version of this license is in
%%    http://www.latex-project.org/lppl.txt
%% and version 1.3 or later is part of all distributions of LaTeX
%% version 2003/12/01 or later.
%%%%%%%%%%%%%%%%

%% If you need to pass whatever options to xcolor
\PassOptionsToPackage{dvipsnames}{xcolor}

%% If you are using \orcid or academicons
%% icons, make sure you have the academicons 
%% option here, and compile with XeLaTeX
%% or LuaLaTeX.
% \documentclass[10pt,a4paper,academicons]{altacv}

%% Use the "normalphoto" option if you want a normal photo instead of cropped to a circle
% \documentclass[10pt,a4paper,normalphoto]{altacv}

\documentclass[10pt,a4paper]{altacv}
%% AltaCV uses the fontawesome and academicon fonts
%% and packages. 
%% See texdoc.net/pkg/fontawecome and http://texdoc.net/pkg/academicons for full list of symbols.
%% 
%% Compile with LuaLaTeX for best results. If you
%% want to use XeLaTeX, you may need to install
%% Academicons.ttf in your operating system's font 
%% folder.


% Change the page layout if you need to
\geometry{left=1cm,right=9cm,marginparwidth=6.8cm,marginparsep=1.2cm,top=1.25cm,bottom=1.25cm,footskip=2\baselineskip}

% Change the font if you want to.

% If using pdflatex:
\usepackage[T1]{fontenc}
\usepackage[utf8]{inputenc}
\usepackage[default]{lato}

% If using xelatex or lualatex:
% \setmainfont{Lato}

% Change the colours if you want to
\definecolor{Mulberry}{HTML}{72243D}
\definecolor{SlateGrey}{HTML}{2E2E2E}
\definecolor{LightGrey}{HTML}{666666}
\definecolor{MyBlue}{HTML}{3388aa}
\definecolor{MyRed}{HTML}{aa3333}
\colorlet{heading}{MyBlue}
\colorlet{accent}{MyRed}
\colorlet{emphasis}{SlateGrey}
\colorlet{body}{LightGrey}

% Change the bullets for itemize and rating marker
% for \cvskill if you want to
\renewcommand{\itemmarker}{{\small\textbullet}}
\renewcommand{\ratingmarker}{\faCircle}
%% sample.bib contains your publications
% \addbibresource{sample.bib}

\usepackage[colorlinks]{hyperref}

\begin{document}

\name{Evan Sitt}
\tagline{Cloud Architect, Backend Developer, \& Support Engineer}
\photo{2.8cm}{White}
\personalinfo{%
  % Not all of these are required!
  % You can add your own with \printinfo{symbol}{detail}
  \email{Sitt.Evan@gmail.com}
  \phone{+1 (325) 939-7642}
  \mailaddress{3429 S. Hermitage Ave, Chicago, IL 60608 United States}
%   \location{San Angelo, Texas, United States}
%   \homepage{http://Yeet.Codes/}
%   \twitter{@twitterhandle}


  \linkedin{linkedin.com/in/evan-sitt/}
  \github{github.com/ParadoxChains}
  %% You MUST add the academicons option to \documentclass, then compile with LuaLaTeX or XeLaTeX, if you want to use \orcid or other academicons commands.
%   \orcid{orcid.org/0000-0000-0000-0000}
}

%% Make the header extend all the way to the right, if you want. 
\begin{fullwidth}
\makecvheader
\end{fullwidth}

%% Depending on your tastes, you may want to make fonts of itemize environments slightly smaller
% \AtBeginEnvironment{itemize}{\small}


%% Provide the file name containing the sidebar contents as an optional parameter to \cvsection.
%% You can always just use \marginpar{...} if you do
%% not need to align the top of the contents to any
%% \cvsection title in the "main" bar.
\cvsection[page1sidebar]{Experience}

\cvevent{Senior Tech Support Engineer}{Microsoft}{June 2021 -- Current}{Dallas, Texas}
\small{
\begin{itemize}
\item Partner with enterprise customers to design, implement, and manage Azure cloud infrastructure tailored to their business needs.

\item Resolve complex technical issues by conducting end-to-end troubleshooting, remediation, and analysis, ensuring minimal downtime and customer impact.

\item Lead cross-team collaboration with product engineering and support teams to address multi-disciplinary challenges and deliver seamless resolutions.

\item Drive product and process improvement by filing bugs, proposing design changes, and validating resolutions to enhance customer experience and operational efficiency.

\item Serve as a trusted advisor, providing guidance on Azure services, architecture best practices, and customer-readiness programs.

\item Contribute to knowledge sharing and readiness initiatives through technical content creation, mentoring, and leading triage sessions.

\end{itemize}
}

\medskip

\divider

\cvevent{Junior Developer}{Farm Fare}{April 2020 -- April 2021}{Cleveland, Ohio}
\small{
\begin{itemize}
\item Review and verify production code written in the Haskell functional programming language before approving pull requests.

\item Build the backend server for handling requests via RESTful API using the Yesod framework written in the Haskell functional programming language.

\item Deploy and maintain back end server on Google Cloud through use of Kubernetes Engine, Cloud Build, and Stackdriver.

\item Ensure Farm Fare's high standard of production code and ease of maintenance via setting up Continuous Integration and Continuous Deployment on code repository. 

\item Configuration of the Odoo business management framework to best fit Farm Fare's objectives and provide detailed end-user documentation.

\end{itemize}
}

\medskip

\divider

\cvevent{Student Developer}{Ericsson Hungary}{March 2019 -- December 2019}{Budapest, Hungary}
\small{
\begin{itemize}
\item Have proper knowledge and skill in coding with Erlang for telecommunication applications.

\item Write functional tests for new functionality developed by the team.

\item Address customer raised Trouble Reports and Issues in a timely manner via debugging and testing.

\item Extend and refactor legacy code for better performance, efficiency, and maintainability.

\end{itemize}
}

\clearpage
\end{document}
